% !TEX encoding = UTF-8 Unicode

%%%%%%%%%%%%%%%%%%%%%%%%%%%%%%%%%%%%%%%%%%%%%%%%%%%%%%%%%%%%%%%
%                                                             %
% El producto de Wallis                                       %
% y $\Gamma(1/2) = \sqrt{\pi}$ son lo mismo                   %
% by                                                          %
% \'Oscar Ciaurri y Juan Luis Varona                          %
%                                                             %
%%%%%%%%%%%%%%%%%%%%%%%%%%%%%%%%%%%%%%%%%%%%%%%%%%%%%%%%%%%%%%%


\documentclass[autocontact]{gaceta}

%---------------------
% La Gaceta se har\'a con esto activado,
% pero no hace falta ahora si no quieres:
\usepackage[T1]{fontenc}
\usepackage{lmodern}

% Puedes usar esto (u otra codificaci\'on) si quieres:
\usepackage[utf8]{inputenc}
\usepackage[spanish]{babel}
%\decimalpoint
%---------------------

%---------------------
% Esto ya lo carga el estilo de La Gaceta:
%\usepackage{amsmath, amsthm, amssymb}
%\usepackage{url} % <-- Para p\'aginas web o similar: \url{...}
%\usepackage{graphicx}
%---------------------

%---------------------
% Carga esto para gr\'aficos si lo necesitas:
% \usepackage{wrapfig}
%---------------------

%---------------------
%\usepackage{tikz}
%---------------------

%---------------------
% <<< ESTO SE AJUSTAR\'A AL EDITAR CADA N\'UMERO DE LA GACETA >>>
%\provisional
\newcounter{pageforminiatura}
\setcounter{pageforminiatura}{566} % Siempre ser\'a p\'agina par
\setcounter{page}{1} % No se usa (dejarlo como 1)
\journame{La Gaceta de la RSME}
\yearofpublication{2018}
\volume{21}
\issuenumber{3}
%---------------------

%---------------------
% Deja esto as\'{\i}:
%\belongstopart{Art\'{\i}culos}
%\belongstopart{Secciones}
%\belongstopart{Noticias de la Sociedad}
%\belongstopart{Actualidad}
\nopart
%---------------------

%---------------------
% Deja esto as\'{\i}:
\supertitle{\large Miniaturas matem\'aticas de \thejourname\\[1mm]
\hfill\footnotesize
Vol.~\thevolume\ (\theyearofpublication),
N\'um.~\theissuenumber, P\'ag.~\thepageforminiatura\theprovisional}
%
\noeditor
%---------------------

\let\contacturl\relax

%---------------------
\title{El producto de Wallis y $\Gamma(1/2) = \sqrt{\pi}$ son lo mismo}
%\author{} % Poner \author{} o no poner nada, da igual.
\contact{\'Oscar Ciaurri, Emilio Fern\'andez y Juan Luis Varona,
Departamento de Matem\'aticas y Computaci\'on,
Universidad de La Rioja}{\texttt{oscar.ciaurri@unirioja.es}, \texttt{eferna35@gmail.com}, \texttt{jvarona@unirioja.es}}
{}
%---------------------

%%%%%%%%%%%%%%%%%%%%%%%%%%%%%%%%%%%%%%%%%%%%%%%%%%%%%%%%%%%%%%%%%%%%%%%
\begin{document}
%%%%%%%%%%%%%%%%%%%%%%%%%%%%%%%%%%%%%%%%%%%%%%%%%%%%%%%%%%%%%%%%%%%%%%%

%%%%%%%%%%%%%%%%%%%%%%%%%%%%%%%%%%%%%%%%%%%%%%%%%%%%%%%%%%%%%%%%%%%%%%%
% Deja esto as\'{\i}:
\previousvspace{-1.8cm}
\makesupertitle
\previousvspace{-0.25cm}
\maketitle
\vspace*{-0.4cm}
\thispagestyle{empty} % <<<<< Importante
%%%%%%%%%%%%%%%%%%%%%%%%%%%%%%%%%%%%%%%%%%%%%%%%%%%%%%%%%%%%%%%%%%%%%%%

Con el objetivo de interpolar entre los términos de la sucesión de factoriales, Euler define la función Gamma en una carta a Goldbach fechada el 13 de octubre de 1729, y lo hace con una expresión que con notación actual se puede escribir como
% Carta a Goldbach fechada el 13 de octubre de 1729
% http://eulerarchive.maa.org/correspondence/correspondents/Goldbach.html
\begin{equation}
\label{eq:Gz}
  \Gamma(z) = \frac{1}{z} \prod_{n=1}^{\infty} \frac{\left(1+\frac{1}{n}\right)^z}{\left(1+\frac{z}{n}\right)}
\end{equation}
(en realidad, Euler usa lo que sería nuestra $\Gamma(z+1)$, y ni el nombre de Gamma ni la notación $\Gamma$ son suyos, sino que los dio Legendre alrededor de 1811). 
% Más adelante, en un artículo publicado en \textit{Commentarii academiae scientiarum Petropolitanae} 
% (vol.~5 (1730/1), 1738, págs.~36--57), da una expresión integral equivalente 
% a la que ahora se toma como definición habitual.
Poco después, en un artículo presentado a la Academia de San Petersburgo el 28 de noviembre de 1729 (publicado en \textit{Comment. acad. sc. Petrop.} \textbf{5} (1730/1), 1738, 36--57), da una expresión integral equivalente a la que ahora se toma como definición habitual.

Usando una versión logarítmica de la identidad $e^{i \pi} = -1$, Euler se da cuenta de que $\text{«}(1/2)!\text{»} = \Gamma(3/2) = \sqrt{\pi}/2$ y lo comunica a Goldbach en la carta antes citada. Además, en el artículo posterior observa ya que esa igualdad (de la que se sigue $\Gamma(1/2) = \sqrt{\pi}$\,) se reduce al producto de Wallis
\[
 \prod_{n=1}^{\infty} \left(\frac{2n}{2n-1} \cdot \frac{2n}{2n+1} \right) =
 \frac{2}{1} \cdot \frac{2}{3} \cdot \frac{4}{3} \cdot \frac{4}{5} \cdot \frac{6}{5}
 \cdot \frac{6}{7} \cdot \frac{8}{7} \cdot \frac{8}{9}  \cdots = \frac{\pi}{2}.
\]
% http://eulerarchive.maa.org//pages/E019.html
Pero este hecho es poco conocido, y el objetivo de esta miniatura es dar una prueba sencilla, con rigor y notación moderna, de que la fórmula $\Gamma(1/2) = \sqrt{\pi}$ y el producto de Wallis son realmente lo mismo.

Para nuestros propósitos, la mejor manera de conectar la función Gamma con un producto infinito es usar \eqref{eq:Gz}, la definición original que dio Euler. Esto implica
\[
\frac{\Gamma(x)\Gamma(y)}{\Gamma(x+y)}
 = \frac{x+y}{xy} \prod_{n=1}^{\infty}
 \frac{\left(1+\frac{x+y}{n}\right)}{\left(1+\frac{x}{n}\right)\left(1+\frac{y}{n}\right)}
 = \frac{x+y}{xy} \prod_{n=1}^{\infty}
 \frac{n(n+x+y)}{(n+x)(n+y)}.
\]
Tomando $x=y=1/2$ en esta expresión y usando $\Gamma(1)=1$, lo que obtenemos es
\begin{align*}
  \Gamma(1/2)^2
  &= \frac{1}{1/4} \prod_{n=1}^{\infty} \frac{n(n+1)}{(n+1/2)(n+1/2)}
  = 4 \prod_{n=1}^{\infty} \frac{2n(2n+2)}{(2n+1)(2n+1)} \\
  &= 4 \cdot \frac{1}{2} \cdot \left( \frac{2}{1} \cdot \frac{2}{3} \cdot \frac{4}{3} \cdot \frac{4}{5}
  \cdot \frac{6}{5} \cdot \frac{6}{7} \cdot \frac{8}{7} \cdot \frac{8}{9}  \cdots\right),
\end{align*}
así que $\Gamma(1/2) = \sqrt{\pi}$ y el producto de Wallis se deducen directamente uno del otro, tal como queríamos comprobar.
% DUDA: "se deducen directamente uno del otro" o "son lo mismo",  
% para que "son lo mismo" no salga tres veces.

\bigskip


%%%%%%%%%%%%%%%%%%%%%%%%%%%%%%%%%%%%%%%%%%%%%%%%%%%%%%%%%%%%%%%%%%%%%%%
\end{document}
%%%%%%%%%%%%%%%%%%%%%%%%%%%%%%%%%%%%%%%%%%%%%%%%%%%%%%%%%%%%%%%%%%%%%%%


